\section{Interpretation and reporting}
\label{ch:conclusion}

A transport improvement first changes the generalized cost of using a link.
Traffic on that link measures the direct saving. At an efficient interior
baseline, the envelope theorem makes this traffic-based statistic the complete
first-order welfare effect: subsequent changes in routes, modes, prices, and
locations have no additional first-order value. This result explains both
the appeal of social-savings calculations and the conditions under which they
are sufficient.

The Extended approach is needed when the baseline contains congestion or
spatial externalities. Route and modal responses then move activity across
wedges, and the primitive policy shock need not pass one for one into the
realized transport cost. Proposition 2 organizes these adjustments as traffic
times pass-through and market-access multipliers. The adjoint calculation
implements the same derivative for the declared set of candidate links without
resolving a separate counterfactual equilibrium for each one. The model
decomposition then asks which modeled effect accounts for the difference
from the Traditional approach.

The economic-geography and urban-commuting applications share this transport
calculation. They differ in the spatial equilibrium that maps transport costs
into production and consumption or into residence and workplace choices.
Consequently, a practitioner must choose the spatial model before
interpreting the welfare result. The same network and traffic data can imply
different welfare derivatives when the underlying economic adjustment margin
changes.

Three checks support different claims. Accounting checks establish that the
input flows, routes, modes, and policy units form a coherent baseline.
Derivative and finite-difference checks establish that the code computes the
stated model. Comparisons with untargeted flows, costs, or outcomes assess
whether that model fits a particular application. Passing one class
of checks does not substitute for passing the others.

The output is a local model-implied benefit from a specified transport-cost
change. It is not a complete project appraisal. Construction and financing
costs, safety, local pollution, displacement, and distributional incidence
must be added separately when they matter for the decision. A reproducible run
therefore reports the policy unit, economic model, parameters, input
transformations, diagnostics, and artifact hashes alongside the welfare
elasticity. Reported in this form, the statistic extends traffic-based
accounting by measuring the spatial-equilibrium correction implied by the
stated model and data. The reference appendices provide the notation-to-code
crosswalk, algorithms, schemas, numerical tests, example results, and
provenance needed to audit that calculation.
